\documentclass[11pt, a4paper]{article}

% ── Packages ──────────────────────────────────────────────────────────
\usepackage[utf8]{inputenc}
\usepackage[T1]{fontenc}
\usepackage{lmodern}
\usepackage[margin=1in]{geometry}
\usepackage{amsmath, amssymb, amsthm}
\usepackage{booktabs}
\usepackage{longtable}
\usepackage{graphicx}
\usepackage{hyperref}
\usepackage{xcolor}
\usepackage{listings}
\usepackage{tcolorbox}
\usepackage{polya}

\hypersetup{
  colorlinks=true,
  linkcolor=polyablue,
  urlcolor=polyablue,
  citecolor=polyablue,
}

% ── Code listing style ────────────────────────────────────────────────
\lstdefinestyle{polya-python}{
  language=Python,
  basicstyle=\ttfamily\small,
  keywordstyle=\color{polyablue}\bfseries,
  commentstyle=\color{polyagray}\itshape,
  stringstyle=\color{polyagreen},
  numbers=left,
  numberstyle=\tiny\color{polyagray},
  numbersep=8pt,
  frame=single,
  framerule=0.4pt,
  rulecolor=\color{polyagray},
  breaklines=true,
  showstringspaces=false,
  tabsize=4,
}
\lstset{style=polya-python}
\title{Debt Payoff Optimization}
\author{uber-polya}
\date{2026-02-26}

\begin{document}

\maketitle
\tableofcontents
\newpage

% ── Phase A: Problem Formulation ─────────────────────────────────────
\section{Problem Formulation}

\begin{polyamodel}

\textbf{Problem Type}: Find \hfill
\textbf{Domain}: Financial Mathematics


\end{polyamodel}

% ── Universe ──────────────────────────────────────────────────────────
\subsection{Universe}

\begin{itemize}
  \item $Problem: \{debts, extra_monthly_payment, total_balance, total_minimums\}$
  \item $Avalanche: \{strategy, total_interest_paid, months_to_payoff, total_paid, is_optimal, is_feasible, algorithm, time_seconds, payoff_order\}$
  \item $Snowball: \{strategy, total_interest_paid, months_to_payoff, total_paid, is_optimal, is_feasible, algorithm, time_seconds\}$
\end{itemize}

% ── Variables ─────────────────────────────────────────────────────────
\subsection{Variables}

\begin{longtable}{lll}
\toprule
\textbf{Symbol} & \textbf{Meaning} & \textbf{Type / Range} \\
\midrule
\endhead
$L$ & loan principal = L & $\mathbb{{R}}_{>0}$ \\
$r$ & monthly interest rate = APR/12 & $\mathbb{{R}}_{>0}$ \\
$n$ & total number of monthly payments & $\mathbb{{Z}}^+$ \\
$PMT$ & monthly payment amount & $\mathbb{{R}}_{>0}$ \\
\bottomrule
\end{longtable}

% ── Structure ─────────────────────────────────────────────────────────
\subsection{Structure}

You have 4 debts with different balances, interest rates, and minimum payments:

| Debt | Balance | APR | Min Payment |
|------|---------|-----|-------------|
| Credit Card | \$6,200 | 22.99\% | \$120/mo |
| Car Loan | \$12,000 | 6.50\% | \$250/mo |
| Student Loan | \$25,000 | 4.50\% | \$280/mo |
| Personal Loan | \$3,500 | 15.00\% | \$75/mo |

**Total debt**: \$46,700. **Total minimums**: \$725/mo. **Extra available**: \$500/mo.

Find the optimal allocation of extra payments to minimize total interest paid
over time. Compare two strategies:

1. **Avalanche** (highest interest rate first) -- mathematically optimal among
   greedy single-target strategies
2. **Snowball** (lowest balance first) -- psychologically motivating but costs
   more in interest

% ── Mapping ───────────────────────────────────────────────────────────
\subsection{Real-World Mapping}

\begin{longtable}{ll}
\toprule
\textbf{Real-World Concept} & \textbf{Mathematical Object} \\
\midrule
\endhead
loan balance & $L — principal$ \\
annual percentage rate & $r — monthly rate = APR/12$ \\
loan term & $n — number of monthly payments$ \\
monthly mortgage payment & $PMT — annuity payment$ \\
\bottomrule
\end{longtable}

% ── Constraints ───────────────────────────────────────────────────────
\subsection{Constraints}

\begin{enumerate}
  \item $PMT = L \cdot \frac{r(1+r)^n}{(1+r)^n - 1}$
  \hfill \textit{(annuity payment formula)}
  \item $B_k = B_{k-1}(1+r) - PMT$
  \hfill \textit{(amortization recurrence)}
  \item $B_n = 0$
  \hfill \textit{(loan fully repaid at maturity)}
  \item $\text{Total Interest} = n \cdot PMT - L$
  \hfill \textit{(total interest = payments minus principal)}
\end{enumerate}

% ── Objective / Claim ─────────────────────────────────────────────────
\subsection{Objective}

\begin{equation}
\min_{\text{option}} \; \sum_{k=1}^{n} PMT_k \; + \; \text{closing costs}
\end{equation}

% ── Approach ──────────────────────────────────────────────────────────
\subsection{Approach}

\begin{description}
  \item[Suggested approach:] unknown
  \item[Complexity class:] 
  \item[Available tools:] unknown
\end{description}
% ── Phase B: Solution ────────────────────────────────────────────────
\section{Solution}

\begin{polyaresult}

\textbf{Answer}: Financial analysis complete


\textbf{Optimal}: No / Unknown \hfill
\textbf{Feasible}: Yes
\textbf{Algorithm}: unknown \quad $$

\textbf{Time}: 0.0s

\textbf{Certificate}: Amortization verified: final balance = \$0.00 for all options.

\end{polyaresult}

% ── Solution Details ──────────────────────────────────────────────────
\subsection{Solution Details}



% ── Verification ──────────────────────────────────────────────────────
\subsection{Verification}

No independent verification checks recorded.

% ── Phase C: Interpretation ──────────────────────────────────────────
\section{Interpretation}

% ── The Question & The Answer ─────────────────────────────────────────
\subsection{The Question}
You have 4 debts with different balances, interest rates, and minimum payments:

| Debt | Balance | APR | Min Payment |
|------|---------|-----|-------------|
| Credit Card | \$6,200 | 22.

\subsection{The Answer}

\begin{polyainsight}[title={Bottom Line}]
Financial analysis complete.
\end{polyainsight}

% ── What This Means ───────────────────────────────────────────────────
\subsection{What This Means}

- Total interest paid under each strategy
- Months to complete payoff under each strategy
- Head-to-head comparison table showing interest savings
- Payoff timeline showing when each debt reaches zero (avalanche)
- Independent verification (5 checks per strategy)

Month-by-month greedy simulation:

1. Accrue interest on all debts (monthly\_rate = APR / 12)
2. Pay minimums on every debt with remaining balance
3. Allocate extra (plus freed minimums from paid-off debts) to the
   highest-priority target debt
4. Cascade surplus -- if the target is paid off mid-allocation, redirect
   remaining funds to the next priority debt
5. Repeat until all balances reach zero

The avalanche strategy prioritizes debts by descending APR, which is provably
optimal because it minimizes the balance-weighted average interest rate at every
step. The snowball strategy prioritizes by ascending balance for faster
psychological wins.

% ── Sensitivity Analysis ──────────────────────────────────────────────

% ── Figures ───────────────────────────────────────────────────────────

% ── Recommendations ───────────────────────────────────────────────────
\subsection{Recommendations}

\begin{enumerate}
  \item Compare total interest paid, not just monthly payments, when evaluating mortgage options.
\end{enumerate}

% ── Limitations ───────────────────────────────────────────────────────
\subsection{Limitations}

\begin{polyawarnbox}
\begin{itemize}
  \item Assumes fixed interest rates; adjustable-rate mortgages introduce variability.
  \item Does not account for tax deductions, insurance, or opportunity cost of down payment.
\end{itemize}
\end{polyawarnbox}

% ── Appendix ─────────────────────────────────────────────────────────

\end{document}